% concatenated from an_upper_bound_for_an_exceptional_automorphism_group.tex
\documentclass[11pt,reqno]{amsart}

\usepackage[T1]{fontenc}
\usepackage{lmodern}
\usepackage{microtype}
\usepackage{amsmath,amssymb,amsthm}
\usepackage{xcolor}
\usepackage[a4paper,margin=1.18in]{geometry}
\usepackage[
  backend=biber,
  style=numeric,
  sorting=none,
  giveninits=true,
  maxbibnames=99,
  doi=true,
  url=false,
  isbn=false
]{biblatex}
\addbibresource{an_upper_bound_for_an_exceptional_automorphism_group.bib}
\usepackage[
  colorlinks=true,
  linkcolor=blue!55!black,
  citecolor=blue!55!black,
  urlcolor=blue!55!black
]{hyperref}
\hypersetup{
  pdftitle={An upper bound for an exceptional automorphism group},
  pdfauthor={Xu Zhuang},
}

\newtheorem{theorem}{Theorem}[section]
\newtheorem{proposition}[theorem]{Proposition}
\newtheorem{lemma}[theorem]{Lemma}
\theoremstyle{remark}
\newtheorem{remark}[theorem]{Remark}

\newcommand{\Aut}{\operatorname{Aut}}
\newcommand{\Gal}{\operatorname{Gal}}
\newcommand{\Fix}{\operatorname{Fix}}
\newcommand{\PGL}{\operatorname{PGL}}
\newcommand{\F}{\mathbf{F}}
\newcommand{\PP}{\mathbf{P}}

\title
{An upper bound for an exceptional automorphism group}

\author{Xu Zhuang}
\address{Department of Mathematics, University of California, Irvine,
  Irvine, CA 92697, USA}
\email{xzhuang8@uci.edu}

\date{\today}



\begin{document}

\begin{abstract}
Let $q=p^h>7$ be odd, put $m=(q+1)/2$, and suppose that $i=(m-2)/2$ satisfies $\gcd(i,m)=\gcd(i+2,m)=1$. For the $\mathbf{F}_{q^2}$-maximal function field $\mathcal{F}_i=\mathbf{F}_{q^2}(x,y)$ defined by $y^m=x^i(x^2+1)$, Peter Beelen, Maria Montanucci, Jonathan Niemann, and Luciane Quoos showed that the geometric automorphism group contains a subgroup of order $4(q+1)$ and conjectured that its order is exactly $4(q+1)$. We prove this equality by establishing the reverse inequality.
\end{abstract}

\maketitle

\section{Introduction}

Maximal function fields over finite fields provide a rich source of curves with many rational points and often with substantial symmetry. Their automorphism groups are useful both as geometric invariants and in the study of isomorphism classes. For a curve given by an explicit equation, one can frequently construct a large subgroup of automorphisms directly. Showing that the constructed subgroup is the full automorphism group is usually a separate problem and may require a geometric description that is intrinsic to the curve.

Let $q=p^h$ be a power of an odd prime and put $m=(q+1)/2$. Beelen, Montanucci, Niemann, and Quoos study the family of $\F_{q^2}$-maximal function fields $\mathcal{F}_i=\F_{q^2}(x,y)$ defined by $y^m=x^i(x^2+1)$,
under the assumptions $\gcd(i,m)=\gcd(i+2,m)=1$. Let $k=\overline{\F}_{q^2}=\overline{\F}_q$ and write $K_i=k\mathcal{F}_i=k(x,y)$ for the extension of constants to $k$. The corresponding smooth projective curve has genus $m-1$; see \cite[Section~3]{BeelenMontanucciNiemannQuoos}. They give a precise description of its geometric automorphism group $\Aut_k(K_i)$ except in the case $q>7$ and $i=(m-2)/2$; see \cite[Theorem~4.9]{BeelenMontanucciNiemannQuoos}. This parameter is fixed by the symmetry $i\mapsto m-2-i$ coming from the reciprocal change of variables in \cite[Lemma~3.2]{BeelenMontanucciNiemannQuoos}, and the exceptional curve consequently has an additional visible symmetry.

In this remaining case, they show that $\Aut_k(K_i)$ contains a subgroup of order $4(q+1)$, which we denote by $H$; see \cite[Theorem~4.6]{BeelenMontanucciNiemannQuoos}. They conjecture that $|\Aut_k(K_i)|=4(q+1)$ \cite[Conjecture~4.11]{BeelenMontanucciNiemannQuoos}. Since $H$ already has this order, the conjecture is equivalent to the assertion that $H$ is the full geometric automorphism group.
\begin{theorem}\label{thm:main}
Let $q=p^h>7$ be odd, let $m=(q+1)/2$, and suppose that $i=(m-2)/2$ satisfies $\gcd(i,m)=\gcd(i+2,m)=1$. If $H$ is the subgroup constructed in \cite[Theorem~4.6]{BeelenMontanucciNiemannQuoos}, then $\Aut_k(K_i)=H$. In particular, $|\Aut_k(K_i)|=4(q+1)$.
\end{theorem}

The lower bound in Theorem~\ref{thm:main} is exactly the result of \cite[Theorem~4.6]{BeelenMontanucciNiemannQuoos}. The content of this paper is the upper bound $|\Aut_k(K_i)|\leq4(q+1)$.

The main idea is to use a quotient different from the original Kummer map $x$. Write $m=2n$, so that the exceptional parameter is $i=n-1$, and consider $u=y^2/x$. The sign involution $\tau:(x,y)\mapsto(x,-y)$ and the reciprocal involution $\rho:(x,y)\mapsto(x^{-1},y/x)$ both fix $u$. We show that they generate the Galois group $V$ of the degree-four extension $K_i/k(u)$. Moreover, $k(u)$ is the unique rational subfield of index four and is preserved by every automorphism of the curve.

It remains to bound the induced action on the $u$-line. The group $V$ has three nonidentity involutions. The sign involution $\tau$ fixes four points of the curve, whereas $\rho$ and $\tau\rho$ each fix $2n$ points. Every automorphism of the curve preserves $k(u)$ and hence conjugates $V$ to itself. Since conjugation preserves the number of fixed points, it must send $\tau$ to itself. The four fixed points of $\tau$ have $u$-values $0$ and $\infty$, so the induced M\"{o}bius transformation preserves the pair $\{0,\infty\}$. It also preserves the remaining branch points, which form the set $\{a\in k:a^{2n}=4\}$. A direct count gives at most $4n$ such transformations and therefore the upper bound $16n=4(q+1)$ for the full automorphism group.


\section{The biquadratic cover and its intermediate quotients}

Since $i=(m-2)/2$ is an integer, $m$ is even. Write $m=2n$, so that $i=n-1$ and $q+1=4n$. The coprimality assumptions imply that $n$ is even: if $n$ were odd, both $n-1$ and $n+1$ would have a common factor $2$ with $2n$. Moreover, $n\geq6$. Indeed, the smaller even values $n=2$ and $n=4$ give respectively $q=7$ and $q=15$, and neither is allowed by the hypotheses.

\begin{remark}
Conversely, when $n$ is even, the equalities $\gcd(n-1,2n)=\gcd(n+1,2n)=1$ hold automatically. Since $q+1=4n$, the exceptional parameters satisfying the coprimality assumptions are therefore exactly the odd prime powers $q>7$ with $q\equiv7\pmod{8}$. The first values are $q=23,31,47,71,79,\ldots$.
\end{remark}

Let $X$ be the smooth projective curve with function field $K=k(X)=K_i=k(x,y)$, where $y^{2n}=x^{n-1}(x^2+1)$.
Choose $\alpha\in k$ with $\alpha^2=-1$. The function $x^{n-1}(x^2+1)$ has zeros of orders $n-1,1,1$ at $x=0,\alpha,-\alpha$, respectively, and a pole of order $n+1$ at $x=\infty$. The coprimality assumptions are $\gcd(n-1,2n)=\gcd(n+1,2n)=1$, so all four orders are coprime to $2n$. Since $p\nmid2n$, the Kummer ramification criterion shows that $x:X\to\PP^1$ is totally ramified exactly above $0,\infty,\alpha,-\alpha$. Let $P_0,P_\infty,P_\alpha,P_{-\alpha}$ denote the unique points of $X$ above them. We use the divisor calculation $(x)=2n(P_0-P_\infty)$ and $(y)=(n-1)P_0+P_\alpha+P_{-\alpha}-(n+1)P_\infty$ from \cite[Equation~(3)]{BeelenMontanucciNiemannQuoos}.
They also show that $g(X)=2n-1$.
Set $u=y^2/x$. The preceding divisor formulas give $(u)=2P_\alpha+2P_{-\alpha}-2P_0-2P_\infty$.
Thus the pole divisor of $u$ has degree four, and hence $[K:k(u)]=4$; see \cite[Theorem~1.4.11]{Stichtenoth}. We will also use the identity $u^n=x+x^{-1}$.

Consider the two automorphisms $\tau(x,y)=(x,-y)$ and $\rho(x,y)=(x^{-1},y/x)$.
The second is the reciprocal isomorphism of \cite[Lemma~3.2]{BeelenMontanucciNiemannQuoos}, specialized to the fixed parameter $i=n-1$. The maps $\tau$ and $\rho$ are commuting involutions, and both fix $u$.

\begin{proposition}\label{prop:cover}
The extension $K/k(u)$ is Galois with group $V=\langle\tau,\rho\rangle\simeq C_2\times C_2$.
\end{proposition}

\begin{proof}
The distinct commuting involutions $\tau$ and $\rho$ generate a subgroup
$V$ of order four. Since both fix $u$, we have
$V\leq\Aut_{k(u)}(K)$. Then $[K:k(u)]=4$ gives
$4=|V|\leq|\Aut_{k(u)}(K)|\leq[K:k(u)]=4$.
Thus $\Aut_{k(u)}(K)=V$, and its order equals the degree of the
extension. Hence $K/k(u)$ is Galois with Galois group $V$.
\end{proof}



We first determine the fixed points of the three nontrivial elements of $V$. Since the characteristic is odd, all three involutions are tame.

\begin{lemma}\label{lem:fixed}
The involution $\tau$ fixes precisely $P_0,P_\infty,P_\alpha,P_{-\alpha}$. The fixed points of $\rho$ are the $2n$ points above $x=1$, and the fixed points of $\tau\rho$ are the $2n$ points above $x=-1$. In particular, $|\Fix(\tau)|=4$ and $|\Fix(\rho)|=|\Fix(\tau\rho)|=2n$.
\end{lemma}

\begin{proof}
The four points $P_0,P_\infty,P_\alpha,P_{-\alpha}$ are the only
ramification points of the Kummer cover $x$, and they are all totally
ramified. Hence every deck transformation fixes them, while the deck
group acts freely elsewhere. Therefore $\tau$ fixes precisely these
four points. Suppose that a point is fixed by $\rho$. Its $x$-coordinate must then be fixed by $x\mapsto x^{-1}$. The points above $0$ and $\infty$ are exchanged, as are $P_\alpha$ and $P_{-\alpha}$, so a fixed point must lie above $x=1$ or $x=-1$. Above $x=1$, the map $\rho$ fixes $y$, and the equation becomes $y^{2n}=2$, which has $2n$ distinct solutions because $p\nmid2n$. Above $x=-1$, the map sends $y$ to $-y$, and no point is fixed because there $y^{2n}=-2\ne0$. The proof for $\tau\rho$ is the same, with $x=1$ and $x=-1$ interchanged.
\end{proof}

\begin{proposition}\label{prop:quotient-genus}
The quotient curve $X/\langle\tau\rangle$ has genus $n-1$, while
$X/\langle\rho\rangle$ and $X/\langle\tau\rho\rangle$ each have genus
$n/2$. In particular, each quotient curve has genus at least three.
\end{proposition}

\begin{proof}
Put $Y=X/\langle\iota\rangle$. The quotient map $X\to Y$ has degree
two, and its ramification points are exactly the fixed points of
$\iota$. Each such point has ramification index two. Since
$\operatorname{char}k\ne2$, this index is prime to the characteristic,
so the map is tame and each ramification point contributes $2-1=1$ to
the ramification term in the Riemann--Hurwitz formula. Consequently, if
$\iota$ has $r$ fixed points, then
$2g(X)-2=2(2g(Y)-2)+r$; see
\cite[Theorem~3.4.13]{Stichtenoth}.

Now $2g(X)-2=4n-4$. For $\iota=\tau$, Lemma~\ref{lem:fixed} gives
$r=4$, and hence $4n-4=2(2g(Y)-2)+4$, so $g(Y)=n-1$.
For $\iota=\rho$ or $\tau\rho$, it gives $r=2n$, and hence
$4n-4=2(2g(Y)-2)+2n$, so $g(Y)=n/2$. Thus each quotient curve has
genus at least three because $n\geq6$.
\end{proof}

\section{The unique rational subfield of index four}

We use Castelnuovo's inequality in the following form. If $F_1,F_2\subset L$, if $L=F_1F_2$, and if $d_j=[L:F_j]$, then
\begin{equation}\label{eq:CS}
  g(L)\leq d_1g(F_1)+d_2g(F_2)+(d_1-1)(d_2-1).
\end{equation}
See \cite[Theorem~3.11.3]{Stichtenoth}.

The first consequence rules out the only low-genus double quotients that could arise in the uniqueness argument.

\begin{proposition}\label{prop:no-low-genus}
There is no degree-two morphism from $X$ to a curve of genus at most one.
\end{proposition}

\begin{proof}
Suppose, to the contrary, that $f:X\to D$ is a degree-two morphism with $g(D)\leq1$, and identify $A=k(D)$ with its image under the pullback $f^*:k(D)\hookrightarrow K$. Then $[K:A]=\deg(f)=2$. Since $\operatorname{char}(k)\ne2$, the extension $K/A$ is separable and, being quadratic, is Galois. Put $B=k(u)$, so $[K:B]=4$. Since $A\subseteq AB\subseteq K$ and $[K:A]=2$, either $AB=A$ or $AB=K$.

If $AB=K$, then $B=k(u)$ is rational, so $g(B)=0$. Applying \eqref{eq:CS} to $A$ and $B$ gives $g(X)\leq2g(D)+4g(B)+(2-1)(4-1)=2g(D)+3\leq5$, contrary to $g(X)=2n-1\geq11$.

It follows that $AB=A$, so $B\subseteq A$. Let $\iota$ be the nontrivial automorphism of the quadratic Galois extension $K/A$. Since $\iota$ fixes $B$, Proposition~\ref{prop:cover} gives $\iota\in V$. Hence $A=K^{\langle\iota\rangle}$ is one of the three intermediate quadratic fields of $K/B$. Proposition~\ref{prop:quotient-genus} shows that each of these fields has genus at least three, a contradiction.
\end{proof}



\begin{proposition}\label{prop:unique}
If $v\in K\setminus k$ satisfies $[K:k(v)]=4$, then $k(v)=k(u)$. Equivalently, $k(u)$ is the unique rational subfield of $K$ having index four.
\end{proposition}

\begin{proof}
Let $B=k(u)$ and $B'=k(v)$, set $L=BB'$, and put $e=[K:L]$. Since $[K:B]=[K:B']=4$, the integer $e$ belongs to $\{1,2,4\}$.

If $e=1$, then $K=BB'$. Applying \eqref{eq:CS} to the two rational fields $B$ and $B'$ gives $g(X)\leq(4-1)^2=9$, contrary to $g(X)\geq11$.

Suppose that $e=2$. The tower law gives $[L:B]=[L:B']=2$, and $L=BB'$. Applying \eqref{eq:CS} inside $L$ gives $g(L)\leq(2-1)^2=1$. Since $[K:L]=2$, this contradicts Proposition~\ref{prop:no-low-genus}.

Therefore $e=4$. Since $B\subseteq L$ and both fields have index four in $K$, one has $L=B$. The same argument gives $L=B'$, and hence $B=B'$.
\end{proof}

\section{The action on the \texorpdfstring{$u$}{u}-line and the upper bound}

Let $A=\Aut_k(K)$. For every $\sigma\in A$, the field $\sigma(k(u))$ is again a rational subfield of index four. Proposition~\ref{prop:unique} therefore implies that $\sigma(k(u))=k(u)$. Thus restriction to $k(u)$ defines a homomorphism $A\to\Aut_k(k(u))\simeq\PGL_2(k)$. If $\Gamma$ is its image, then
\begin{equation*}
1\longrightarrow V\longrightarrow A\longrightarrow\Gamma\longrightarrow1
\end{equation*}
is exact.

We next describe the branch points of the biquadratic cover
$u:X\to\PP^1$. The four fixed points of $\tau$ map to $0$ or $\infty$
by the divisor of $u$. At a fixed point of $\rho$ one has $x=1$, so
$u^n=x+x^{-1}=2$; at a fixed point of $\tau\rho$ one has $x=-1$, so
$u^n=-2$. Conversely, the points above $x=1$ and $x=-1$ map onto all
roots of these two equations. Hence the three branch subsets associated
with $\tau$, $\rho$, and $\tau\rho$ are respectively $\{0,\infty\}$,
$\{a\in k:a^n=2\}$, and $\{a\in k:a^n=-2\}$. They are pairwise
disjoint, and the full branch locus is $\mathcal{B}=\{0,\infty\}\cup R$,
where $R=\{a\in k:a^{2n}=4\}$.

The subgroup $V$ is normal in $A$, since it is the kernel of the restriction map. For $\sigma\in A$, the conjugate $\sigma\tau\sigma^{-1}$ belongs to $V$ and has the same number of fixed points as $\tau$. By Lemma~\ref{lem:fixed}, $\tau$ is the unique nontrivial element of $V$ with four fixed points, because $2n>4$. Thus $\sigma\tau\sigma^{-1}=\tau$ for every $\sigma\in A$. It follows that $\sigma$ preserves $\Fix(\tau)$, and consequently the induced element of $\Gamma$ preserves its image $\{0,\infty\}$ on the $u$-line. Conjugation by $\sigma$ fixes $\tau$ and therefore permutes the remaining two nontrivial elements $\rho$ and $\tau\rho$ of $V$. Since $\sigma(\Fix(\iota))=\Fix(\sigma\iota\sigma^{-1})$ for every $\iota\in V$, the automorphism $\sigma$ preserves $\Fix(\rho)\cup\Fix(\tau\rho)$. Since $u(\Fix(\rho)\cup\Fix(\tau\rho))=R$, the induced element of $\Gamma$ preserves $R$.

A M\"{o}bius transformation preserving the unordered pair $\{0,\infty\}$ has the form $t\mapsto ct$ or $t\mapsto c/t$, with $c\in k^\times$. The first form preserves $R$ precisely when $c^{2n}=1$, while the second preserves $R$ precisely when $c^{2n}=16$. Since $p\nmid2n$, each equation has exactly $2n$ solutions in $k$. Therefore $|\Gamma|\leq4n$, and hence $|A|=|V||\Gamma|\leq4\cdot4n=16n=4(q+1)$.

\begin{proof}[Proof of Theorem~\ref{thm:main}]
The preceding argument gives $|\Aut_k(K_i)|\leq4(q+1)$. By \cite[Theorem~4.6]{BeelenMontanucciNiemannQuoos}, the group $\Aut_k(K_i)$ contains the subgroup $H$ of order $4(q+1)$. Hence equality holds and $\Aut_k(K_i)=H$.
\end{proof}

\enlargethispage{3\baselineskip}
\section*{Acknowledgments}
The author thanks Yuxiang Yao for helpful discussions.

\section*{Statement on AI-assisted preparation}
During the preparation of this manuscript, the author used OpenAI's GPT-5.6 Sol
language model through Codex for exploratory proof development, literature
organization, and language editing. In particular, the suggestion to organize
the upper-bound argument around the common invariant $u=y^2/x$ arose during
that interaction. The author independently checked every argument 
and assumes full responsibility for the content of the manuscript.

\printbibliography[title={References}]

\end{document}
